
\chapter{相关说明}

\section{字体兼容问题}

\subsection{Windows 本地编译}

无需任何配置，开箱即用。

\subsection{非Windows平台}

系统字体与 Windows 存在差异，可能影响排版细节。即使不配置字体，模板也能正常编译通过，仅在笔画细节上与标准 Windows 环境略有差异。

\subsubsection{推荐解决方案}

\begin{enumerate}
    \item 使用 TeXPage（推荐）：平台内置了标准的 Windows 字体库，无需手动配置。
    \item 手动配置字体：将宋体、黑体、仿宋、楷体等字体文件从 Windows 系统拷贝至项目根目录的 fonts/ 文件夹。
\end{enumerate}

\begin{table}[htbp]
    \centering
    \caption{各平台中文字体对照表}
    \label{tab:fonts}
    \begin{tabular}{llll}
        \toprule
        类别 & Windows (中易出品) & Mac (华文出品) & Linux (开源/Fandol) \\
        \midrule
        宋体 & SimSun (simsun.ttc) & STSong & FandolSong \\
        黑体 & SimHei (simhei.ttf) & STHeiti & FandolHei \\
        楷体 & SimKai (simkai.ttf) & STKaiti & FandolKai \\
        仿宋 & SimFang (simfang.ttf) & STFangsong & FandolFang \\
        版权 & 商业版权 & 商业版权 & 免费开源 \\
        \bottomrule
    \end{tabular}
\end{table}

\section{项目文件结构}

项目主要包含以下文件和目录：

\begin{itemize}
    \item \texttt{nkuthesis.cls}：模板核心文件
    \item \texttt{main.tex}：硕博主文档（修改个人信息和内容）
    \item \texttt{main.bib}：参考文献数据库
    \item \texttt{main-bachelor.tex}：本科主文档
    \item \texttt{data/}：章节内容文件夹
    \item \texttt{assets/}：资源文件目录
    \item \texttt{figures/}：图片文件夹
    \item \texttt{fonts/}：字体文件目录（可选）
\end{itemize}

\section{快速开始}

\subsection{编译器选择}

本模板必须使用 XeLaTeX 引擎编译，参考文献需配合 biber 处理。

\subsubsection{推荐在线平台}

\begin{itemize}
    \item TeXPage：国内访问更稳定，原生支持中文模板
    \item Overleaf：选择 XeLaTeX 作为编译器
\end{itemize}

\subsubsection{本地编译环境}

\begin{itemize}
    \item Windows：安装 TeX Live 或 MiKTeX
    \item Mac：安装 MacTeX
    \item 推荐编辑器：VS Code + LaTeX Workshop 插件
\end{itemize}

\subsection{编译方式}

推荐编译顺序：xelatex main.tex → biber main → xelatex main.tex → xelatex main.tex

\subsection{学位类型配置}

\begin{itemize}
    \item 本科使用 main-bachelor.tex
    \item 硕博在 main.tex 中指定：\texttt{\textbackslash documentclass[master, arabicnumber]\{nkuthesis\}}
    \item 可选学位：master（硕士）、doctor（博士）
    \item 章节编号：arabicnumber（1.1样式）或 chinesenumber（第一章样式）
\end{itemize}

\subsection{个人信息填写}

所有个人信息通过 \texttt{\textbackslash nkusetup\{\}} 统一设置，主要包括论文题目、作者信息、指导教师、学科专业等。

\section{本章小结}

本章介绍了南开大学 \LaTeX 模板的使用环境、文件结构和基本配置方法。通过合理配置，用户可以快速开始学位论文的撰写工作。
